\documentclass[a4paper,11pt]{letter}
\usepackage[utf8]{inputenc}
\usepackage[T1]{fontenc}
\usepackage{tgheros}
\renewcommand{\familydefault}{\sfdefault}
\usepackage[english]{babel}
\usepackage[colorlinks=true, urlcolor=blue]{hyperref}
\usepackage{geometry}
\geometry{left=1in, right=1in, top=1in, bottom=1in}

\name{Baha Khammari}
\address{D\"usseldorf, Germany \\
    \href{mailto:b.e.khammari@gmail.com}{b.e.khammari@gmail.com} \\
    +49\,1520\,9016956}

\begin{document}

\begin{letter}{%
    % --- ADAPT PER APPLICATION ---
    [Professor / Selection Committee] \\
    [Department of Economics] \\
    [University] \\
    [Address]
}

\opening{Dear [Professor / Members of the Selection Committee],}

I am writing to apply for the [PhD position / doctoral programme in \ldots] at [University / Department]. I am completing my M.Sc.\ in Economics at the University of Cologne (expected September 2026, current GPA: 1.6/1.0) and will be available from October 2026. My research interests lie in applied microeconomics, causal inference, and development economics, with a methodological focus on difference-in-differences estimation and heterogeneous treatment effects.

% --- RESEARCH BACKGROUND ---

My master's thesis, graded 1.0/1.0 and supervised by Anna Person, M.Sc., provides the first causal analysis of Brazil's ProUni scholarship programme on formal-sector wages at the microregion level. I constructed a panel of 558 regions over 2005--2019 from administrative microdata ($\sim$5\,million RAIS observations) and applied the Callaway \& Sant'Anna (2021) heterogeneity-robust staggered DiD estimator with doubly-robust inference. The key finding---a non-monotonic dose--response with significant wage gains in low-exposure regions that vanish at higher intensities---was entirely masked by conventional two-way fixed effects, illustrating the empirical relevance of recent methodological advances in treatment-effect heterogeneity.

This work taught me to manage large-scale administrative datasets (SQL/BigQuery), implement frontier econometric methods in R, and think carefully about the gap between average effects and distributional consequences. A separate seminar paper with Prof.\ Dr.\ Michael Krause on housing and the macroeconomy, analysing subprime mortgage credit through the Mian \& Sufi framework, broadened my perspective on how micro-level frictions propagate through aggregate channels.

% --- FIT AND MOTIVATION ---
% ADAPT: Explain why this specific programme/supervisor.
% Reference their work where it connects to your agenda.

I am particularly drawn to [programme / your research group] because [specific reason: a paper they published, a methodological focus, a dataset, a policy lab, a research cluster]. My proposed doctoral research---a comparative causal evaluation of Brazil's three major higher education access policies (ProUni, FIES, Quota Law)---would complement [their existing work on \ldots]. The project is empirically feasible using publicly available Brazilian administrative data with which I have extensive experience, and the methodological toolkit aligns with [the department's strengths in \ldots].

% --- PROFESSIONAL BACKGROUND ---

Beyond research, two years as an intern and working student at PwC Germany (Assurance Solutions) gave me experience with large-scale data analysis under time pressure, rigorous documentation standards, and professional communication with senior stakeholders. An Erasmus+ semester at RSM Erasmus University (Rotterdam) exposed me to an international academic environment and broadened my analytical perspective through coursework in asset pricing and derivatives valuation.

% --- CLOSING ---

I would welcome the opportunity to discuss how my research agenda aligns with [the programme's / your group's] direction. My academic CV, a detailed research proposal, and a thesis abstract are enclosed. I have asked Anna Person, M.Sc.\ (master's thesis supervisor), Univ.-Prof.\ Dr.\ Andreas Schabert (bachelor's thesis supervisor), and Prof.\ Dr.\ Michael Krause (seminar instructor) to serve as referees; their contact details are available in my CV.

Thank you for considering my application.

\closing{Yours sincerely,}

\end{letter}

\end{document}
